% !TEX TS-program = LuaLaTeX
% !TEX encoding = UTF-8 Unicode

%% ==========================================================================
%% Vorlage für eine Bachelorarbeit an der IU Internationalen Hochschule
%% Konform mit den offiziellen Richtlinien (Stand: 01.10.2025)
%% ==========================================================================

\documentclass[11pt, a4paper, oneside, ngerman]{article}

%% --------------------------------------------------------------------------
%% Sprach- und Texteinstellungen
%% --------------------------------------------------------------------------
\usepackage[ngerman]{babel}                 % Deutsche Trennmuster und Begriffe
\usepackage[autostyle=true, german=quotes]{csquotes} % Korrekte deutsche Anführungszeichen

%% --------------------------------------------------------------------------
%% Schriftarten & Typografie (LuaLaTeX / fontspec)
%% --------------------------------------------------------------------------
\usepackage{fontspec}
\usepackage{microtype}                      % Randausgleich und optimierter Zeilenumbruch

% Schriftart Arial gemäß IU-Leitfaden (mit Ausweichschriftarten für Linux/CI)
\IfFontExistsTF{Arial}{
    \setmainfont{Arial}
}{
    \IfFontExistsTF{Liberation Sans}{
        \setmainfont{Liberation Sans}
    }{
        \setmainfont{TeX Gyre Heros}
    }
}

%% --------------------------------------------------------------------------
%% Seitenränder & Layout
%% --------------------------------------------------------------------------
% IU-Richtlinie: Randbreiten betragen exakt 2,00 cm für alle Ränder
\usepackage[
    left=2.0cm,
    right=2.0cm,
    top=2.0cm,
    bottom=2.0cm,
    headsep=0.8cm,
    footskip=1.0cm
]{geometry}

% Zeilenabstand: 1,5-zeilig im Fließtext (IU-Vorgabe)
\usepackage{setspace}
\setstretch{1.5}

% Textabsätze: 6 Pt. Abstand nach Absatz, kein Erstzeileneinzug (IU-Vorgabe)
\usepackage[skip=6pt]{parskip}

% Fußnoten: 10 Pt. im Blocksatz (IU-Vorgabe)
\renewcommand{\footnotesize}{\fontsize{10pt}{13pt}\selectfont}

%% --------------------------------------------------------------------------
%% Überschriftenformatierung (Exakt gemäß IU-Tabelle 1 & Anhang B)
%% --------------------------------------------------------------------------
% 1. Ebene: 16 Pt. fett, linksbündig | Abstand: 12 Pt. vor, 12 Pt. nach
% 2. Ebene: 14 Pt. fett, linksbündig | Abstand: 12 Pt. vor, 6 Pt. nach
% 3. Ebene: 11 Pt. fett, linksbündig | Abstand: 12 Pt. vor, 6 Pt. nach
\usepackage{titlesec}
\titleformat{\section}{\normalfont\fontsize{16pt}{20pt}\bfseries}{\thesection}{1em}{}
\titlespacing*{\section}{0pt}{12pt}{12pt}

\titleformat{\subsection}{\normalfont\fontsize{14pt}{18pt}\bfseries}{\thesubsection}{1em}{}
\titlespacing*{\subsection}{0pt}{12pt}{6pt}

\titleformat{\subsubsection}{\normalfont\fontsize{11pt}{15pt}\bfseries}{\thesubsubsection}{1em}{}
\titlespacing*{\subsubsection}{0pt}{12pt}{6pt}

%% --------------------------------------------------------------------------
%% Inhaltsverzeichnis (Formatierung gemäß IU-Richtlinien)
%% --------------------------------------------------------------------------
% Nur Überschriften der 1. Ebene fett, weitere Ebenen normal
% Überschriften und Verzeichnisse linksbündig, durchgehende Punktlinie
\makeatletter
\renewcommand*\l@section[2]{%
  \ifnum \c@tocdepth >\z@
    \addpenalty\@secpenalty
    \addvspace{0.8em \@plus\p@}%
    \setlength\@tempdima{1.5em}%
    \begingroup
      \parindent \z@ \rightskip \@pnumwidth
      \parfillskip -\@pnumwidth
      \leavevmode \bfseries
      \advance\leftskip\@tempdima
      \hskip -\leftskip
      #1\nobreak\normalfont\leaders\hbox{$\m@th
        \mkern \@dotsep mu\hbox{.}\mkern \@dotsep
        mu$}\hfill\nobreak\hb@xt@\@pnumwidth{\hss \bfseries #2}\par
    \endgroup
  \fi}
\makeatother

%% --------------------------------------------------------------------------
%% Kopf- und Fußzeilen (Seitennummerierung zentriert am Seitenende)
%% --------------------------------------------------------------------------
\usepackage{fancyhdr}
\pagestyle{fancy}
\fancyhf{}
\cfoot{\thepage}                            % Zentrierte Seitenzahl am Seitenende
\renewcommand{\headrulewidth}{0pt}
\renewcommand{\footrulewidth}{0pt}
\fancypagestyle{plain}{
    \fancyhf{}
    \cfoot{\thepage}
    \renewcommand{\headrulewidth}{0pt}
}

%% --------------------------------------------------------------------------
%% Grafiken, Tabellen & Quellenangaben
%% --------------------------------------------------------------------------
\usepackage{graphicx}

% Gleitumgebungen standardmäßig zentrieren
\makeatletter
\g@addto@macro\@floatboxreset\centering
\makeatother

% Einheitliche Quellenangabe unter Abbildungen und Tabellen (10 Pt. gemäß IU-Richtlinien)
\newcommand{\quelle}[1]{%
    \vspace{4pt}%
    \par
    \begingroup
    \footnotesize
    \noindent Quelle: #1\par
    \endgroup
}

%% --------------------------------------------------------------------------
%% Hyperlinks & PDF-Metadaten
%% --------------------------------------------------------------------------
\usepackage[hidelinks]{hyperref}
\urlstyle{same}                             % URLs in Fließtext-Schriftart

\hypersetup{
    pdfauthor={<Vorname Nachname>},
    pdftitle={<Titel der Bachelorarbeit>},
    pdfsubject={Bachelorarbeit -- IU Internationale Hochschule},
    pdfkeywords={IU, Internationale Hochschule, Bachelorarbeit},
    pdfcreator={LuaLaTeX},
    pdfproducer={LuaLaTeX}
}

%% --------------------------------------------------------------------------
%% Hilfspakete
%% --------------------------------------------------------------------------
\usepackage{lipsum}
\usepackage[nohyperlinks, printonlyused]{acronym}

%% --------------------------------------------------------------------------
%% Quellcode-Formatierung (listings)
%% --------------------------------------------------------------------------
\usepackage{listings}
\usepackage{xcolor}

\definecolor{codegreen}{rgb}{0,0.6,0}
\definecolor{codegray}{rgb}{0.5,0.5,0.5}
\definecolor{codepurple}{rgb}{0.58,0,0.82}
\definecolor{backcolour}{rgb}{0.97,0.97,0.97}

\lstset{
    backgroundcolor=\color{backcolour},
    commentstyle=\color{codegreen},
    keywordstyle=\color{blue}\bfseries,
    numberstyle=\tiny\color{codegray},
    stringstyle=\color{codepurple},
    basicstyle=\small\ttfamily,
    breaklines=true,
    breakatwhitespace=true,
    showstringspaces=false,
    columns=flexible,
    numbers=none,
    tabsize=2,
    frame=tb
}

%% --------------------------------------------------------------------------
%% Literaturverzeichnis (BibLaTeX nach APA 7th Edition)
%% --------------------------------------------------------------------------
\usepackage[
    backend=biber,
    style=apa,
    pagetracker,
    apamaxprtauth=20
]{biblatex}

% Komma vor dem letzten Autor bei deutschem 'und' entfernen (APA-Konvention im Deutschen)
\makeatletter
\renewcommand*{\apablx@ifrevnameappcomma}{\@secondoftwo}
\makeatother

\DefineBibliographyExtras{ngerman}{%
    \renewcommand*{\finalandcomma}{}%
}

\addbibresource{references.bib}

%% ==========================================================================
%% DOKUMENTENSTRUKTUR GEMÄSS IU-RICHTLINIEN
%% ==========================================================================
\begin{document}

% --------------------------------------------------------------------------
% 1. Titelblatt (Zählt als Seite I, Nummerierung wird nicht gedruckt)
% --------------------------------------------------------------------------
\hypersetup{pageanchor=false}
\pagenumbering{Roman}
%% Title Page
\def\usesf{}
\let\usesf\sffamily % Uncomment this line for the default TeX font

\newsavebox{\Tutor}
\savebox{\Tutor}{<Name of the Tutor>}


\setlength{\unitlength}{1pt}

\begin{titlepage}
\begin{center}
\includegraphics[width=.32\paperwidth]{logos/iu_Logo_EN.jpg}

\hbox{}
\vfill
{\usesf}
{\huge\bfseries My Case Study \par}
\vskip 1.8cm
Case Study\\[2mm]
\vskip 1cm

{\large\bfseries Markus Haug (<Student ID>)\\}
% Neu: Anschrift hinzugefügt entsprechend der IU Muster-Titelseite
<Musterstraße 50>\\
<12345 Musterstadt>\\
\vskip 1.2cm
Studiengang: Business Informatics (B.Sc.)\\
% in the Module\\
ABCDF01 - Course Code\\
\today % Abgabedatum
\vskip 3cm
\begin{tabular}{p{3cm}l}
Tutor: & \usebox{\Tutor} \\
\end{tabular}
\vfill
\end{center}

\end{titlepage}
%% End of Title Page


%%% Local Variables:
%%% mode: latex
%%% TeX-master: "thesis"
%%% End:
\hypersetup{pageanchor=true}

% --------------------------------------------------------------------------
% Vorspann: Römische Paginierung beginnt auf Seite II
% --------------------------------------------------------------------------
\setcounter{page}{2}

% 2. Danksagung (freiwillig bei Abschlussarbeiten; bei Bedarf einkommentieren)
% %% ==========================================================================
%% Danksagung (freiwillig bei Abschlussarbeiten gemäß IU-Richtlinien)
%% ==========================================================================

\phantomsection
\section*{Danksagung}
\addcontentsline{toc}{section}{Danksagung}

An dieser Stelle möchte ich mich bei allen Personen bedanken, die mich während der Ausarbeitung dieser Bachelorarbeit unterstützt und motiviert haben.

Ein besonderer Dank gilt meiner Erstprüferin / meinem Erstprüfer für die fachliche Begleitung, die wertvollen Anregungen und die konstruktive Betreuung während des gesamten Arbeitsprozesses.

Ebenso bedanke ich mich herzlich bei meiner Familie und meinen Freunden für ihren Rückhalt und ihr Verständnis während dieser intensiven Studienphase.

% \newpage

% 3. Gender-Disclaimer (optional gemäß Anhang A; bei Bedarf einkommentieren)
% %% ==========================================================================
%% Gender-Disclaimer (optional gemäß IU-Richtlinien Anhang A)
%% ==========================================================================

\phantomsection
\section*{Gender-Disclaimer}
\addcontentsline{toc}{section}{Gender-Disclaimer}

\textbf{Hinweis zur geschlechtergerechten Sprache:}

Aus Gründen der besseren Lesbarkeit wird in diesem Dokument überwiegend das generische Maskulinum verwendet. Selbstverständlich sind damit Personen aller Geschlechter gleichermaßen gemeint.

% \newpage

% 4. Abstract (Pflicht bei Abschlussarbeiten laut Thesis-Handbuch)
\phantomsection
\section*{Abstract}
\addcontentsline{toc}{section}{Abstract}
\noindent
\lipsum[1] % TODO: Hier den tatsächlichen Abstract einfügen

\newpage

% 5. Inhaltsverzeichnis (erscheint nicht als Eintrag im Inhaltsverzeichnis)
\tableofcontents
\newpage

% 6. Abbildungsverzeichnis (Pflicht ab 3 Abbildungen)
\phantomsection
\addcontentsline{toc}{section}{\listfigurename}
\listoffigures
\newpage

% 7. Tabellenverzeichnis (Pflicht ab 3 Tabellen)
\phantomsection
\addcontentsline{toc}{section}{\listtablename}
\listoftables
\newpage

% 8. Abkürzungsverzeichnis
\phantomsection
\section*{Abkürzungsverzeichnis}
\addcontentsline{toc}{section}{Abkürzungsverzeichnis}
\begin{acronym}[IU]
    \acro{iu}[IU]{IU Internationale Hochschule}
\end{acronym}

\newpage

% --------------------------------------------------------------------------
% 9. Hauptteil (Textteil: Arabische Paginierung beginnend bei Seite 1)
% --------------------------------------------------------------------------
\pagenumbering{arabic}

%% ==========================================================================
%% Einleitung (Beginn des bewerteten Textteils, Seite 1 arabisch)
%% ==========================================================================

\section{Einleitung}
\label{sec:einleitung}

% --------------------------------------------------------------------------
% Leitfaden für die Einleitung einer Bachelorarbeit:
% 1. Ausgangslage, Problemstellung und Relevanz des Themas
% 2. Zielsetzung der Arbeit und forschungsleitende Fragestellung(en)
% 3. Methodischer Ansatz und theoretischer Bezugsrahmen
% 4. Aufbau und Gang der Untersuchung
% --------------------------------------------------------------------------

In der heutigen wissenschaftlichen und unternehmerischen Praxis gewinnt die fundierte Analyse und methodische Aufarbeitung komplexer Sachverhalte zunehmend an Bedeutung. Die vorliegende Bachelorarbeit widmet sich der zentralen Fragestellung, wie moderne Software- und Architekturkonzepte zielgerichtet bewertet und eingesetzt werden können.

\lipsum[1-2]
\lipsum[1-2]   % Kapitel 1: Einleitung
\newpage

%% ==========================================================================
%% Hauptteil der Bachelorarbeit
%% ==========================================================================

\section{Hauptteil}
\label{sec:hauptteil}

\lipsum[1]

\subsection{Theoretische Grundlagen}
\label{subsec:theorie}

Wissenschaftliche Arbeiten erfordern eine fundierte Auseinandersetzung mit dem aktuellen Forschungsstand. Nach \textcite{blinowski_monolithic_2022} lassen sich verschiedene Architekturansätze bezüglich ihrer Performanz und Skalierbarkeit differenziert evaluieren. Auch grundlegende methodische Ansätze können über direkte und indirekte Zitate eingebunden werden \parencite[S.~42]{nussbaum}.

\subsubsection{Konzepte und Begriffsabgrenzung}
\label{subsubsec:konzepte}

Gemäß den IU-Richtlinien können Arbeiten in maximal drei Gliederungsstufen unterteilt werden. Hierbei erfolgt die Formatierung der Überschriften der 3.~Stufe in 11\,Pt. fett. 

\lipsum[2]

\subsection{Methodik und empirische Untersuchung}
\label{subsec:methodik}

Die empirische Erhebung erfolgte anhand leitfadengestützter Experteninterviews. Der vollständige Erhebungsbogen ist im Anhang dokumentiert (siehe Anhang~A).

In Abbildung~\ref{fig:logo_iu} ist beispielhaft das Logo der \ac{iu} dargestellt. Abbildungen und Tabellen sollten stets im Fließtext referenziert und mit einer eindeutigen Quellenangabe in Schriftgröße 10\,Pt. versehen werden.

\begin{figure}[htbp]
    \centering
    \includegraphics[width=0.45\textwidth]{logos/iu_Logo_EN.jpg}
    \caption[Logo der \acs{iu}]{Logo der \ac{iu}}
    \label{fig:logo_iu}
    \quelle{In Anlehnung an \url{https://www.iu.de}}
\end{figure}

Zur Veranschaulichung von strukturierten Daten dient Tabelle~\ref{tab:beispiel}. Laut IU-Richtlinien ist ein Tabellenverzeichnis ab drei verwendeten Tabellen verpflichtend.

\begin{table}[htbp]
    \centering
    \caption{Übersicht exemplarischer Bewertungskriterien}
    \label{tab:beispiel}
    \vspace{4pt}
    \begin{tabular}{llr}
        \hline
        \textbf{Kriterium} & \textbf{Beschreibung} & \textbf{Gewichtung} \\
        \hline
        Inhaltliche Tiefe & Systematische Bearbeitung der Forschungsfrage & 40\,\% \\
        Wissenschaftlichkeit & Literaturarbeit nach APA-Richtlinien & 30\,\% \\
        Struktur \& Formalia & Einhaltung des IU-Leitfadens und Layouts & 30\,\% \\
        \hline
    \end{tabular}
    \quelle{Eigene Darstellung}
\end{table}

\subsection{Quellcode-Darstellung}
\label{subsec:code}

Quellcode kann übersichtlich über das \texttt{listings}-Paket eingebunden werden:

\begin{lstlisting}[language=Java, caption={Beispielhafte Java-Methode}, label={lst:java_beispiel}]
public class Application {
    public static void main(String[] args) {
        System.out.println("Willkommen an der IU!");
    }
}
\end{lstlisting}

\lipsum[7]
\lipsum[7-8]       % Kapitel 2: Hauptteil
\newpage

%% ==========================================================================
%% Fazit und Ausblick
%% ==========================================================================

\section{Fazit und Ausblick}
\label{sec:fazit}

% --------------------------------------------------------------------------
% Leitfaden für das Fazit einer Bachelorarbeit:
% 1. Zusammenfassende Beantwortung der Forschungsfragen
% 2. Diskussion der zentralen Erkenntnisse und Implikationen für Theorie & Praxis
% 3. Kritische Reflexion der Methodik und Limitationen der Arbeit
% 4. Ausblick auf zukünftigen Forschungs- und Entwicklungsbedarf
% --------------------------------------------------------------------------

Die vorliegende Bachelorarbeit hat sich mit der systematischen Untersuchung der Forschungsfrage auseinandergesetzt. Zusammenfassend konnte gezeigt werden, dass die gewählte Methodik belastbare Ergebnisse liefert und praktische Implikationen für das Fachgebiet aufzeigt.

\lipsum[4-5]

     % Kapitel 3: Fazit und Ausblick
\newpage

% --------------------------------------------------------------------------
% 10. Literaturverzeichnis (fortlaufende arabische Seitennummerierung)
% --------------------------------------------------------------------------
\printbibliography[heading=bibintoc, title={Literaturverzeichnis}]

% --------------------------------------------------------------------------
% 11. Anhang mit Anhangsverzeichnis (fortlaufende arabische Seitennummerierung)
% --------------------------------------------------------------------------
\newpage
%% ==========================================================================
%% Anhang & Anhangsverzeichnis (Exakt gemäß IU-Richtlinien)
%% ==========================================================================

% --------------------------------------------------------------------------
% Anhangsverzeichnis (verpflichtend bei mehreren Anhängen)
% Laut IU-Richtlinien: Unmittelbar vor dem eigentlichen Anhang eingefügt,
% enthält die Bezeichnungen der Anhänge ohne Seitenzahlen.
% Erscheint als 'Anhangsverzeichnis' im Hauptinhaltsverzeichnis.
% --------------------------------------------------------------------------
\phantomsection
\section*{Anhangsverzeichnis}
\addcontentsline{toc}{section}{Anhangsverzeichnis}

\vspace{0.8cm}

\noindent
\textbf{Anhang A:} Interviewleitfaden und Erhebungsinstrumente\\[0.4cm]
\textbf{Anhang B:} Ergänzende statistische Auswertungen und Datensätze\\[0.4cm]
\textbf{Anhang C:} Quellcode und technische Konfigurationen

\newpage

% --------------------------------------------------------------------------
% Anhang A
% --------------------------------------------------------------------------
\phantomsection
\section*{Anhang A: Interviewleitfaden und Erhebungsinstrumente}
\label{anhang:A}

Dieser Anhang enthält den für die empirische Erhebung verwendeten Interviewleitfaden. Im Textteil der Bachelorarbeit muss mindestens einmal auf diesen Anhang verwiesen werden (z.\,B. \enquote{siehe Anhang~A}).

\lipsum[1-2]

\newpage

% --------------------------------------------------------------------------
% Anhang B
% --------------------------------------------------------------------------
\phantomsection
\section*{Anhang B: Ergänzende statistische Auswertungen und Datensätze}
\label{anhang:B}

Hier können umfangreiche Tabellen, Fragebögen oder statistische Auswertungen dokumentiert werden, die den fortlaufenden Lesefluss im Hauptteil stören würden.

\lipsum[3]
\lipsum[1-2]

% --------------------------------------------------------------------------
% 12. Ehrenwörtliche Erklärung (Pflichtbestandteil an der IU)
% --------------------------------------------------------------------------
\newpage
%% ==========================================
%% Ehrenwörtliche Erklärung
%% ==========================================

\phantomsection
\section*{Ehrenwörtliche Erklärung}
\addcontentsline{toc}{section}{Ehrenwörtliche Erklärung}

Ich versichere hiermit, dass ich die vorliegende Arbeit selbstständig verfasst und keine anderen als die angegebenen Quellen und Hilfsmittel benutzt habe. Alle Ausführungen, die wörtlich oder sinngemäß übernommen wurden, sind als solche kenntlich gemacht.

Diese Arbeit hat in gleicher oder ähnlicher Form noch keiner anderen Prüfungsbehörde vorgelegen und wurde auch noch nicht veröffentlicht.

\vspace{2.5cm}

\noindent
\begin{tabular}{@{}p{6.5cm}p{1.5cm}p{6.5cm}@{}}
\dotfill & & \dotfill \\
Ort, Datum & & Eigenhändige Unterschrift \\
\end{tabular}


\end{document}
\end{document}